\documentclass[12pt,letterpaper]{article}
\usepackage[utf8]{inputenc}
\usepackage[T1]{fontenc}
\usepackage{amsmath,amssymb}
\usepackage{graphicx}
\usepackage{booktabs}
\usepackage{hyperref}
\usepackage[margin=1in]{geometry}
\usepackage{natbib}
\bibliographystyle{apalike}

\title{Parallel Prefix Scan for Atmospheric Radiative Transfer:\\
A 50-Year Gap Between Theory and Practice}

\author{
  [Author Name]\\
  \textit{Independent Researcher}\\
  \texttt{[email]}\\
  \\
  Code: \url{https://github.com/consigcody94/parallel-prefix-rt}
}

\date{March 2026}

\begin{document}
\maketitle

\begin{abstract}
Radiation schemes are the most computationally expensive physics component in numerical weather prediction models. The vertical transport computation---the ``adding method''---has been inherently sequential in every weather model, requiring $O(N_{\text{layers}})$ steps that cannot be parallelized on GPUs. We present the first application of parallel prefix scan to atmospheric radiative transfer by reformulating the adding method's recurrence as a suffix scan of $2\times 2$ M\"{o}bius transformation matrices. Matrix multiplication is associative, enabling the Blelloch parallel scan algorithm to reduce sequential depth from $O(N)$ to $O(\log N)$. On an NVIDIA RTX 3060 with 128 atmospheric layers, the parallel scan achieves a \textbf{4.73$\times$ speedup} with maximum relative error $7.76 \times 10^{-7}$ and zero accuracy failures across 132,096 test points. The mathematical basis---that radiative transfer layer operators form a semigroup under the Redheffer star product---was established by Grant and Hunt (1969), and the identical parallel scan technique has been deployed in recurrent neural networks (Martin \& Cundy, 2018) and state space models (Gu \& Dao, 2023, Mamba). Yet this connection to atmospheric radiative transfer remained unexploited for over 50 years. We additionally quantify single-precision speedup (7.98$\times$ on GPU) and demonstrate a fast exponential approximation (2.55$\times$ on CPU, sub-ULP accuracy). All code is open source.
\end{abstract}

\section{Introduction}

Atmospheric radiation schemes compute the transfer of shortwave and longwave radiation through the atmosphere, accounting for absorption and scattering by gases, clouds, and aerosols. These calculations consume 30--50\% of total runtime in numerical weather prediction (NWP) and climate models \citep{pincus2019,ukkonen2024}.

The most recent state-of-the-art result is that of \citet{ukkonen2024}, who achieved a 12$\times$ speedup on ECMWF's ecRad radiation scheme through spectral reduction, code restructuring, and single-precision arithmetic. However, one fundamental bottleneck remained: the vertical transport computation. The adding method \citep{shonk2008}, which propagates diffuse radiation through the atmospheric column, requires two sequential passes---bottom-up for albedo and top-down for fluxes---creating an $O(N_{\text{layers}})$ sequential dependency.

\subsection{The Key Insight}

The adding method's recurrence for layer albedo:
\begin{equation}
\alpha_i = R_i + \frac{T_i^2 \, \alpha_{i+1}}{1 - R_i \, \alpha_{i+1}}
\label{eq:adding}
\end{equation}
is a M\"{o}bius (linear fractional) transformation $f(x) = (ax+b)/(cx+d)$, representable as a $2\times 2$ matrix:
\begin{equation}
\mathbf{M}_i = \begin{pmatrix} T_i^2 - R_i^2 & R_i \\ -R_i & 1 \end{pmatrix}
\label{eq:mobius}
\end{equation}
Composition of M\"{o}bius transformations corresponds to matrix multiplication, which is \textbf{associative}---the only property required for parallel prefix scan \citep{blelloch1990}.

\subsection{The 50-Year Gap}

\citet{grant1969} proved in 1969 that radiative transfer layer operators form a semigroup under the Redheffer star product \citep{redheffer1959}, establishing associativity. \citet{blelloch1990} published the general parallel prefix scan algorithm in 1990. \citet{martin2018} parallelized mathematically identical linear recurrences in RNNs on GPUs in 2018, achieving 9$\times$ speedup. \citet{gu2023} deployed the same associative scan in the Mamba state space model in 2023. Yet no prior work connected these results to atmospheric radiative transfer.

\section{Method}

Each atmospheric layer $i$ is converted to the $2\times 2$ matrix (\ref{eq:mobius}). The recurrence (\ref{eq:adding}) requires the suffix product:
\begin{equation}
\mathbf{S}_k = \mathbf{M}_k \cdot \mathbf{M}_{k+1} \cdots \mathbf{M}_{N-1}
\end{equation}
followed by application to the surface boundary condition:
\begin{equation}
\alpha_k = \frac{S_k^{(a)} \cdot \alpha_{\text{sfc}} + S_k^{(b)}}{S_k^{(c)} \cdot \alpha_{\text{sfc}} + S_k^{(d)}}
\end{equation}

The suffix product is computed via right-to-left Hillis-Steele inclusive scan in GPU shared memory, requiring $\lceil \log_2 N \rceil$ parallel steps. For $N=128$ layers: 7 steps instead of 128.

\section{Results}

All benchmarks performed on an NVIDIA GeForce RTX 3060 (12 GB, compute capability 8.6, CUDA 13.2).

\subsection{Parallel Prefix Scan}

\begin{table}[h]
\centering
\begin{tabular}{lccr}
\toprule
Method & Time (ms) & Speedup & Max relative error \\
\midrule
Sequential (1 thread/column) & 0.203 & 1.00$\times$ & Reference \\
\textbf{Parallel scan} (1 block/column) & \textbf{0.043} & \textbf{4.73$\times$} & $7.76 \times 10^{-7}$ \\
\bottomrule
\end{tabular}
\caption{Parallel prefix scan vs.\ sequential adding method. 1024 columns, 128 layers. Zero accuracy failures across 132,096 test points.}
\label{tab:scan}
\end{table}

\subsection{Precision Optimization}

\begin{table}[h]
\centering
\begin{tabular}{lccc}
\toprule
Precision & Time (ms) & Bandwidth (GB/s) & Speedup \\
\midrule
FP64 & 18.86 & 28.5 & 1.00$\times$ \\
\textbf{FP32} & \textbf{2.36} & \textbf{113.6} & \textbf{7.98$\times$} \\
\bottomrule
\end{tabular}
\caption{Transmissivity computation ($33.5 \times 10^6$ elements). The 7.98$\times$ speedup reflects the RTX 3060's 1:32 FP64:FP32 throughput ratio.}
\label{tab:precision}
\end{table}

\subsection{Fast Exponential (CPU)}

\begin{table}[h]
\centering
\begin{tabular}{lccc}
\toprule
Method & Time (ms) & Speedup & Max relative error \\
\midrule
Intrinsic exp() & 4.375 & 1.00$\times$ & Reference \\
\textbf{Fast exp} (12th-order Horner) & \textbf{1.719} & \textbf{2.55$\times$} & $3.15 \times 10^{-16}$ \\
\bottomrule
\end{tabular}
\caption{Transmissivity array computation on CPU (gfortran 15.2, -O3 -march=native, double precision). Sub-ULP accuracy confirmed across $10^7$ test points.}
\label{tab:fastexp}
\end{table}

\section{Discussion}

The theoretical maximum speedup from the parallel scan is $N/\log_2 N = 128/7 \approx 18\times$. The achieved 4.73$\times$ reflects synchronization overhead, shared memory bank conflicts, and the 2$\times$ work factor of Hillis-Steele. A work-efficient Blelloch scan or warp-level primitives could improve the ratio.

The parallel scan is applicable to any model using the adding/doubling method: NOAA GFS (RTE+RRTMGP), ECMWF IFS (ecRad), DWD ICON, NCAR MPAS, and DOE E3SM. The M\"{o}bius transformation property is inherent to the physics.

On GPU, hardware exp() is already optimal ($\sim$1.0$\times$ from our polynomial), so the GPU strategy should focus on precision reduction and algorithmic parallelization. On CPU, fast transcendental functions provide the main opportunity.

\section{Conclusion}

We have presented the first parallel prefix scan for atmospheric radiative transfer, achieving 4.73$\times$ GPU speedup with perfect accuracy. The mathematical foundations existed for 50+ years, the parallel algorithm for 35 years, and successful GPU implementations for identical structures for 8 years---yet the connection was never made. Code is available at \url{https://github.com/consigcody94/parallel-prefix-rt}.

\begin{thebibliography}{99}
\bibitem[Blelloch(1990)]{blelloch1990} Blelloch, G.E. (1990). Prefix Sums and Their Applications. CMU-CS-90-190.
\bibitem[Grant \& Hunt(1969)]{grant1969} Grant, I.P. and Hunt, G.E. (1969). Discrete space theory of radiative transfer. \textit{Proc.\ R.\ Soc.\ A}, 313, 183--197.
\bibitem[Gu \& Dao(2023)]{gu2023} Gu, A. and Dao, T. (2023). Mamba: Linear-Time Sequence Modeling with Selective State Spaces. \textit{arXiv:2312.00752}.
\bibitem[Martin \& Cundy(2018)]{martin2018} Martin, E. and Cundy, C. (2018). Parallelizing Linear Recurrent Neural Nets Over Sequence Length. \textit{ICLR}.
\bibitem[Pincus \& Mlawer(2019)]{pincus2019} Pincus, R. and Mlawer, E.J. (2019). Balancing Accuracy, Efficiency, and Flexibility in Radiation Calculations for Dynamical Models. \textit{JAMES}, 11, 3074--3089.
\bibitem[Redheffer(1959)]{redheffer1959} Redheffer, R.M. (1959). Inequalities for a matrix Riccati equation. \textit{J.\ Math.\ Mech.}, 8, 349--367.
\bibitem[Shonk \& Hogan(2008)]{shonk2008} Shonk, J.K.P. and Hogan, R.J. (2008). Tripleclouds. \textit{J.\ Climate}, 21, 2352--2370.
\bibitem[Ukkonen \& Hogan(2024)]{ukkonen2024} Ukkonen, P. and Hogan, R.J. (2024). Twelve Times Faster yet Accurate. \textit{JAMES}, 16. doi:10.1029/2023MS003932.
\end{thebibliography}

\end{document}
